% !TEX root = ../algebra_02.tex

\section{Свободные группы}

\subsection{Определение и основные свойства}

\begin{Definition}{Свободная группа}{}
	Пусть $F$ — группа, $f_1, \dots, f_n$ — её элементы. 
	
	Группа $F$ называется свободной со свободными образующими $f_1, \dots, f_n$, если выполняется следующее универсальное свойство: для любой группы $G$ и любых элементов $g_1, \dots, g_n \in G$ существует единственный гомоморфизм $\varphi: F \to G$ такой, что $\varphi(f_i) = g_i$ для $i = 1, \dots, n$.
\end{Definition}

\begin{Example}{Свободная группа с одним образующим}{}
	Группа целых чисел $\mathbb{Z}$ является свободной группой со свободным образующим $1$. 
	
	Действительно, для любого $g \in G$ существует гомоморфизм $\varphi: \mathbb{Z} \to G$, заданный правилом $k \mapsto g^k$, при котором $1 \mapsto g$.
\end{Example}

\begin{Proposition}{Единственность свободной группы}{}
	Пусть $F$ и $F'$ — свободные группы со свободными образующими $f_1, \dots, f_n$ и $f_1', \dots, f_n'$ соответственно. Тогда существует единственный изоморфизм $\varepsilon: F \to F'$ такой, что $\varepsilon(f_i) = f_i'$ для $i = 1, \dots, n$.
\end{Proposition}

\begin{proof}
	По определению свободной группы существует единственный гомоморфизм $\varepsilon: F \to F'$ с требуемым свойством $\varepsilon(f_i) = f_i'$. 
	
	Аналогично, в обратную сторону существует единственный гомоморфизм $\varepsilon': F' \to F$ такой, что $\varepsilon'(f_i') = f_i$ для $i = 1, \dots, n$.
	
	Рассмотрим композицию $\varepsilon' \circ \varepsilon: F \to F$. Для образующих имеем:
	\[ (\varepsilon' \circ \varepsilon)(f_i) = f_i, \quad i = 1, \dots, n \]
	Тождественное отображение $id_F: F \to F$ также удовлетворяет условию $id_F(f_i) = f_i$. Из единственности гомоморфизма следует, что $\varepsilon' \circ \varepsilon = id_F$.
	
	Аналогично показывается, что $\varepsilon \circ \varepsilon' = id_{F'}$. Поскольку гомоморфизм $\varepsilon$ обратим, он является изоморфизмом.
\end{proof}

\begin{Proposition}{Представление конечно порожденной группы}{}
	Пусть группа $G = \langle g_1, \dots, g_n \rangle$ конечно порождена, а $F$ — свободная группа с $n$ свободными образующими. Тогда существует нормальная подгруппа $N \triangleleft F$ такая, что $G \cong F/N$.
\end{Proposition}

\begin{proof}
	Пусть $f_1, \dots, f_n$ — свободные образующие группы $F$. По определению свободной группы существует гомоморфизм $\varphi: F \to G$, переводящий $f_i \mapsto g_i$.
	
	Так как $g_1, \dots, g_n \in \operatorname{Im} \varphi$, а эти элементы порождают $G$, получаем, что $\langle g_1, \dots, g_n \rangle \subset \operatorname{Im} \varphi$, следовательно, $\operatorname{Im} \varphi = G$.
	
	По теореме о гомоморфизме:
	\[ \operatorname{Im} \varphi \cong F / \ker \varphi \]
	Поскольку $\operatorname{Im} \varphi = G$, достаточно положить $N = \ker \varphi$.
\end{proof}

\begin{Example}{Конечно порожденные группы}{}
	Примеры систем образующих для известных групп:
	\begin{itemize}
		\item $\mathbb{Z} \times \mathbb{Z} = \langle (1,0), (0,1) \rangle$
		\item $S_n = \langle (1 \dots n), (12) \rangle$. Заметим, что $(12)^{(1 \dots n)} = (23)$.
		\item $D_n = \langle R_1, S_0 \rangle$
	\end{itemize}
\end{Example}

\subsection{Конструкция свободной группы}

Пусть $n \in \mathbb{N}$. Зададим алфавит $A = \{x_1, x_1', \dots, x_n, x_n'\}$, где символ $x_i'$ играет роль формального обратного элемента для $x_i$.

Определим множество слов $W$:
\[ W = \{a_1 \dots a_k \mid k \ge 0, a_i \in A, i = 1, \dots, k\} \]
Введем операцию конкатенации (умножения) слов $W \times W \to W$:
\[ (a_1 \dots a_k, b_1 \dots b_l) \mapsto a_1 \dots a_k b_1 \dots b_l \]

Будем говорить, что слово $\widetilde{w}$ получается из слова $w$ \textbf{вставкой}, если $w = w_0 w_1$ и $\widetilde{w} = w_0 x_i x_i' w_1$ или $\widetilde{w} = w_0 x_i' x_i w_1$ для некоторых $w_0, w_1 \in W$.

Слово $\widetilde{w}$ получается из $w$ \textbf{сокращением}, если $w$ получается из $\widetilde{w}$ вставкой.

Введем отношение $\sim$: $w \sim \widetilde{w}$, если $\widetilde{w}$ получается из $w$ конечным числом вставок или сокращений. Это отношение является отношением эквивалентности.

Определим множество $F_n$ как фактормножество $W$ по отношению эквивалентности $\sim$:
\[ F_n := W / \sim \]
Обозначим через $[w]$ класс эквивалентности слова $w$. Введем умножение на $F_n$:
\[ F_n \times F_n \to F_n, \quad ([w], [w']) \mapsto [ww'] \]

\begin{Lemma}{Корректность умножения}{}
	Операция умножения на $F_n$ задана корректно, то есть результат не зависит от выбора представителей классов эквивалентности.
\end{Lemma}

\begin{proof}
	Пусть $w \sim \widetilde{w}$ и $w' \sim \widetilde{w}'$. Тогда выполнена цепочка эквивалентностей:
	\[ ww' \sim \widetilde{w}w' \sim \widetilde{w}\widetilde{w}' \implies ww' \sim \widetilde{w}\widetilde{w}' \]
	Следовательно, классы совпадают: $[ww'] = [\widetilde{w}\widetilde{w}']$.
\end{proof}

\begin{Theorem}{Конструкция свободной группы}{}
	Множество $F_n = W / \sim$ является свободной группой с $n$ свободными образующими $f_1, \dots, f_n$, где $f_i := [x_i]$ для $i = 1, \dots, n$.
\end{Theorem}

\begin{proof}
	Сначала покажем, что $F_n$ образует группу:
	\begin{itemize}
		\item Ассоциативность операции очевидна из ассоциативности конкатенации слов.
		\item Нейтральным элементом является класс пустого слова $[]$.
		\item Проверим существование обратных элементов. Заметим, что:
		\[ [x_i] \cdot [x_i'] = [x_i x_i'] = [] \]
		\[ [x_i'] \cdot [x_i] = [x_i' x_i] = [] \]
		В общем случае, пусть $g \in F_n$. Представим его как $g = [a_1 \dots a_m] = [a_1] \dots [a_m]$, где $a_i \in A$. Тогда обратным элементом к $g$ будет класс $[a_m]^{-1} \dots [a_1]^{-1}$, где под $[x_i]^{-1}$ понимается $[x_i']$, а под $[x_i']^{-1}$ понимается $[x_i]$. В этом случае:
		\[ g \cdot \left([a_m]^{-1} \dots [a_1]^{-1}\right) = [] \quad \text{и} \quad \left([a_m]^{-1} \dots [a_1]^{-1}\right) \cdot g = [] \]
	\end{itemize}
	
	Теперь проверим универсальное свойство. Пусть задана группа $G$ и элементы $g_1, \dots, g_n \in G$.
	Определим отображение $\Phi: W \to G$ на образующих:
	\[ \Phi(x_i) = g_i, \quad \Phi(x_i') = g_i^{-1} \]
	И продолжим его на произвольные слова:
	\[ \Phi(a_1 \dots a_m) = \Phi(a_1) \dots \Phi(a_m) \]
	
	Покажем, что если $w \sim w'$, то $\Phi(w) = \Phi(w')$. Достаточно проверить это для одной операции вставки или сокращения. Пусть $\widetilde{w} = v_0 x_i x_i' v_1$, тогда:
	\[ \Phi(v_0 x_i x_i' v_1) = \Phi(v_0) g_i g_i^{-1} \Phi(v_1) = \Phi(v_0)\Phi(v_1) = \Phi(v_0 v_1) \]
	
	Следовательно, отображение корректно спускается на фактормножество, то есть существует отображение $\varphi: F_n \to G$ такое, что $\varphi([w]) = \Phi(w)$.
	Проверим, что $\varphi$ — гомоморфизм:
	\[ \varphi([w][w']) = \varphi([ww']) = \Phi(ww') = \Phi(w)\Phi(w') = \varphi([w])\varphi([w']) \]
	
	Докажем единственность. Пусть существует другой гомоморфизм $\varphi': F_n \to G$, удовлетворяющий условию $\varphi'(f_i) = g_i$. Тогда:
	\[ \varphi'([x_i]) = g_i = \varphi([x_i]) \]
	\[ \varphi'([x_i']) = g_i^{-1} = \varphi([x_i']) \]
	Так как значения на всех образующих и их обратных совпадают, то для любого $g \in F_n$ выполнено $\varphi'(g) = \varphi(g)$. Универсальное свойство доказано.
\end{proof}

\begin{Remark}{Порождение свободной группы}{}
	Из конструкции видно, что группа порождается своими образующими: $F_n = \langle f_1, \dots, f_n \rangle$.
	
	Следовательно, если $F$ — произвольная свободная группа со свободными образующими $f_1^\circ, \dots, f_n^\circ$, то она также порождается ими: $F = \langle f_1^\circ, \dots, f_n^\circ \rangle$.
\end{Remark}

\subsection{Нормальное замыкание и несократимые слова}

\begin{Definition}{Нормальное замыкание}{}
	Пусть $G$ — группа, $M \subset G$. Нормальным замыканием множества $M$ (или нормальной подгруппой, порожденной $M$) называется:
	\[ \langle\langle M \rangle\rangle := \bigcap_{\substack{M \subset H \\ H \triangleleft G}} H \]
	Очевидно, что по построению $\langle\langle M \rangle\rangle \triangleleft G$.
\end{Definition}

\begin{Definition}{Несократимое слово}{}
	Слово $w \in W$ называется несократимым, если в нем нельзя произвести ни одного сокращения (то есть оно не содержит подслов вида $x_i x_i'$ или $x_i' x_i$).
\end{Definition}

\subsection{Несократимые слова}

\begin{Proposition}{О несократимых словах}{}
	Каждый класс эквивалентности (элемент $F_n = W / \sim$) содержит ровно одно несократимое слово.
\end{Proposition}

\begin{proof}
	\textbf{Существование:} очевидно, что в каждом классе эквивалентности есть хотя бы одно несократимое слово, так как длина слова конечна, и процесс последовательных сокращений завершается за конечное число шагов.
	
	\textbf{Единственность:} пусть $w \sim w'$, где $w$ и $w'$ — несократимые слова. Предположим от противного, что $w \neq w'$.
	По определению отношения эквивалентности существует конечная последовательность преобразований:
	\[ \lambda: w = w_0 \mapsto w_1 \mapsto \dots \mapsto w_n = w' \]
	где каждый переход $w_i \mapsto w_{i+1}$ является либо вставкой, либо сокращением подслова $x_i x_i^{-1}$.
	
	Определим длину последовательности преобразований $\lambda$ как сумму длин всех участвующих в ней слов:
	\[ l(\lambda) = l(w_0) + l(w_1) + \dots + l(w_n) \]
	Среди всех возможных последовательностей преобразований, переводящих $w$ в $w'$, выберем ту, для которой значение $l(\lambda)$ минимально.
	
	Далее в выбранной последовательности выберем слово $w_j$ максимальной длины. Если таких слов несколько, выберем слово с наименьшим индексом $j \ge 1$. Поскольку начальное слово $w_0 = w$ несократимо, первый шаг $w_0 \mapsto w_1$ обязательно является вставкой, следовательно, локальный максимум длины существует и $j \ge 1$.
	
	Рассмотрим переходы в окрестности максимума: $w_{j-1} \mapsto w_j$ (обязательно вставка) и $w_j \mapsto w_{j+1}$ (обязательно сокращение). Из комбинаторики слов следует, что можно произвести оба сокращения из $w_j$ (те, которые вели бы обратно в $w_{j-1}$ и в $w_{j+1}$) независимым образом. Это означает, что существует слово $w_j'$, к которому можно прийти из $w_{j-1}$ сокращением, и из которого можно прийти в $w_{j+1}$ вставкой. При этом $l(w_j') = l(w_j) - 4$.
	
	Тогда мы можем заменить фрагмент пути $w_{j-1} \mapsto w_j \mapsto w_{j+1}$ на $w_{j-1} \mapsto w_j' \mapsto w_{j+1}$, что строго уменьшит общую сумму длин слов $l(\lambda)$. Это противоречит выбору последовательности $\lambda$ с минимальным значением $l(\lambda)$. Следовательно, $w = w'$.
\end{proof}

\begin{Remark}{}{}
	Ввиду доказанного утверждения, можно считать, что свободная группа $F_n$ состоит в точности из несократимых слов в алфавите $A$. При этом образующий $x_i$ отождествляется с $f_i$, а $x_i'$ — с $f_i^{-1}$.
\end{Remark}

\subsection{Копредставления групп}

\begin{Definition}{Задание группы образующими и соотношениями}{}
	Запись вида:
	\[ \langle f_1, \dots, f_n \mid r_1, \dots, r_m \rangle \]
	где $r_1, \dots, r_m \in F_n$, обозначает факторгруппу:
	\[ F_n / \langle\langle r_1, \dots, r_m \rangle\rangle \]
	Данная конструкция называется группой, заданной образующими $f_1, \dots, f_n$ и определяющими соотношениями $r_1, \dots, r_m$.
\end{Definition}

\begin{Example}{Копредставление группы диэдра $D_5$}{}
	Покажем, что $D_5 \cong \langle S, T \mid S^2, T^5, (ST)^2 \rangle$.
	
	Пусть $F_2$ — свободная группа со свободными образующими $S, T$. Зададим гомоморфизм $\varphi: F_2 \to D_5$, действующий на образующие следующим образом:
	\[ S \mapsto S_0 \text{ (симметрия)}, \quad T \mapsto R_1 \text{ (поворот)} \]
	В группе $D_5$ элементы $S_0$ и $R_1$ удовлетворяют соотношениям $S_0^2 = e$, $R_1^5 = e$ и $(S_0 R_1)^2 = e$. Следовательно:
	\[ S^2 \in \ker \varphi, \quad T^5 \in \ker \varphi, \quad (ST)^2 \in \ker \varphi \]
	Определим нормальную подгруппу $N := \langle\langle S^2, T^5, STST \rangle\rangle$. Из включений выше очевидно, что $N \subset \ker \varphi$.
	
	Рассмотрим факторгруппу $F_2 / N = \langle s, t \rangle$, где $s = SN$, $t = TN$. По определению факторгруппы по нормальному замыканию, в $F_2/N$ справедливы тождества:
	\[ s^2 = e, \quad t^5 = e, \quad stst = e \]
	Из последнего равенства выразим $st$:
	\[ st = t^{-1}s^{-1} = t^4 s \]
	Используя это соотношение (позволяющее перебрасывать $s$ вправо от $t$), любой элемент $g \in F_2 / N$ можно привести к виду $t^k s^l$:
	\[ g = s^{i_1} t^{j_1} s^{i_2} t^{j_2} \dots = t^k s^l \]
	где показатели степеней ограничены: $0 \le k \le 4$ и $0 \le l \le 1$. Следовательно, порядок факторгруппы ограничен сверху: $|F_2 / N| \le 10$.
	
	По теореме о гомоморфизме, так как $N \subset \ker \varphi$, существует индуцированный гомоморфизм $\widetilde{\varphi}: F_2 / N \to D_5$.
	Поскольку $\langle S_0, R_1 \rangle = D_5$, исходный гомоморфизм $\varphi$ является сюръективным, а значит, и $\widetilde{\varphi}$ сюръективен.
	Мы имеем сюръективное отображение группы, в которой не более 10 элементов, на группу из ровно 10 элементов ($|D_5| = 10$). Это возможно только в том случае, если $|F_2 / N| = 10$ и $\widetilde{\varphi}$ является биекцией.
	Таким образом, $\widetilde{\varphi}$ — изоморфизм.
	
	\textbf{Упражнение:} Задать образующими и соотношениями группу $(\mathbb{Z}/6\mathbb{Z}) \times (\mathbb{Z}/10\mathbb{Z})$.
\end{Example}

\begin{Remark}{Табличное копредставление конечной группы}{}
	Любую конечную группу можно задать образующими и соотношениями.
	Пусть $|G| = n$, а её элементы занумерованы как $G = \{g_1, \dots, g_n\}$. Очевидно, что $G = \langle g_1, \dots, g_n \rangle$.
	Рассмотрим свободную группу $F_n$ со свободными образующими $f_1, \dots, f_n$. Каждому равенству $g_i g_j = g_k$, выполняющемуся в группе $G$, сопоставим соотношение $f_i f_j f_k^{-1}$.
	Тогда группу $G$ можно представить как:
	\[ G \cong \langle f_1, \dots, f_n \mid \{ f_i f_j f_k^{-1} \}_{i,j,k=1}^n \text{ для всех } i,j,k \text{ таких, что } g_i g_j = g_k \text{ в } G \rangle \]
\end{Remark}